\section*{Hoofdstuk \ref{ch:g12:innate-immunity} --- Aangeboren afweer}

\begin{solution}{exo:g12:innate-immunity:1}
Aangeboren: vanaf de geboorte aanwezig, werkt binnen enkele uren, herkent
klassen van indringers aan gedeelde patronen, reageert elke keer gelijk.
Verworven: heeft dagen nodig, herkent elke indringer afzonderlijk,
onthoudt hem (en bestaat alleen bij \omterm{def:g10:common-ancestry:bodyplan}{gewervelden}).
\end{solution}

\begin{solution}{exo:g12:innate-immunity:2}
Roodheid en warmte: verwijde vaten die meer bloed aanvoeren. Zwelling:
lekke vaatwanden die plasma het weefsel in laten. Pijn: zenuwuiteinden die
door de \omterm{prop:g12:innate-immunity:inflammation}{mediatoren} gevoeliger zijn gemaakt.
\end{solution}

\begin{solution}{exo:g12:innate-immunity:3}
Moleculaire patronen die klassen van microben delen (wandbrokstukken,
\omterm{def:g10:chemistry-of-life:families}{nucleïnezuren} van virussen) en moleculen die door beschadigde \omterm{def:g10:cells-common-unit:cell}{cellen}
worden afgegeven. Zij geven \omterm{prop:g12:innate-immunity:inflammation}{mediatoren} af: histamine, prostaglandinen,
cytokinen.
\end{solution}

\begin{solution}{exo:g12:innate-immunity:4}
Contact en binding door ontvangers; opslokken door het membraan;
vertering in een blaasje met \omterm{def:g11:enzymes-and-phenotype:enzyme}{enzymen} en bijtende stoffen; brokstukken van
\omterm{def:g10:chemistry-of-life:families}{eiwitten} getoond op het oppervlak.
\end{solution}

\begin{solution}{exo:g12:innate-immunity:5}
Zij blokkeren het \omterm{def:g11:enzymes-and-phenotype:enzyme}{enzym} dat prostaglandinen maakt, waardoor zwelling, pijn
en koorts afnemen. Zij nemen de oorzaak niet weg en doden geen microbe.
\end{solution}

\begin{solution}{exo:g12:innate-immunity:6}
De aangeboren rond een halve dag; de verworven rond dag 9; de krommen
kruisen elkaar rond dag 5--6.
\end{solution}

\begin{solution}{exo:g12:innate-immunity:7}
De aanleiding is een moleculair patroon dat de wachtcellen herkennen, en
geen levende indringer. Steriel zout water draagt geen patroon en geeft
geen schade: hoogstens een lichte reactie op de naald.
\end{solution}

\begin{solution}{exo:g12:innate-immunity:8}
Dode neutrofielen, dode en levende bacteriën, verteerd puin en plasma. De
neutrofielen komen over enkele uren aan en sterven over een dag; pus is
wat zich ophoopt zodra zij gevochten hebben.
\end{solution}

\begin{solution}{exo:g12:innate-immunity:9}
Hooikoorts is histamine die mestcellen tegen stuifmeel afgeven: histamine
blokkeren neemt de reactie weg. Bij een infectie stuurt histamine alleen
de eerste minuten; de latere zwelling, pijn en oproeping draaien op
prostaglandinen en cytokinen, waar het antihistaminicum niets aan doet.
\end{solution}

\begin{solution}{exo:g12:innate-immunity:10}
Op 12 uur 100\% tegen 55\%; op 36 uur 50\% tegen 30\%. Het middel verlaagt
de zwelling op elk tijdstip maar de kromme zakt nog altijd rond 72 uur tot
nul: het verkort de reactie niet.
\end{solution}

\begin{solution}{exo:g12:innate-immunity:11}
Gescheurde vezels van een band en verpletterde \omterm{def:g10:cells-common-unit:cell}{cellen} geven de moleculen
van beschadigde \omterm{def:g10:cells-common-unit:cell}{cellen} af, die de ontvangers van de wachtcellen herkennen
zoals zij een microbe zouden herkennen: schade alleen is al een
aanleiding.
\end{solution}

\begin{solution}{exo:g12:innate-immunity:12}
Zonder neutrofielen zijn de eerste dagen verloren en overweldigen gewone
bacteriën het lichaam: de \omterm{def:g12:innate-immunity:immunity}{aangeboren afweer} is wat meteen de linie houdt.
Zonder \omterm{def:g12:innate-immunity:immunity}{verworven afweer} ruimt de aangeboren de eerste infecties op, maar
aanhoudende of herhaalde infecties, die eigen antistoffen en geheugen
vergen, winnen op den duur: de \omterm{def:g12:innate-immunity:immunity}{verworven afweer} maakt af wat de aangeboren
in bedwang houdt.
\end{solution}

\begin{solution}{exo:g12:innate-immunity:13}
Een abces is een ontsteking die levende bacteriën bevecht: haar dempen
vertraagt het opruimen en riskeert verspreiding. Gewrichtsontsteking is
een ontsteking zonder enige microbe, gericht op het \omterm{def:g10:muscles-and-joints:joint}{gewricht} zelf: daar is
de reactie de ziekte, en haar verminderen is de behandeling.
\end{solution}

\begin{solution}{exo:g12:innate-immunity:14}
De ontvangers van de \omterm{def:g12:innate-immunity:immunity}{aangeboren afweer} liggen vast in de \omterm{def:g10:universal-dna:gene}{genen} en herkennen
patronen die hele klassen delen; een tweede ontmoeting verandert daar
niets aan. De \omterm{def:g12:innate-immunity:immunity}{verworven afweer} bouwt ontvangers die op elke indringer
afzonderlijk passen en houdt de \omterm{def:g10:cells-common-unit:cell}{cellen} die hen maakten: zij herkent het
individu, en bewaart die herkenning.
\end{solution}

\begin{solution}{exo:g12:innate-immunity:15}
Een ontsteking is de genezing wanneer zij op een echte indringer of een
echt letsel is gericht, in verhouding staat, en binnen enkele dagen
voorbij is: zij ruimt de plek op en zet het herstel in gang. Zij is de
ziekte wanneer zij op iets onschuldigs is gericht (stuifmeel, de eigen
\omterm{def:g10:muscles-and-joints:joint}{gewrichten} van het lichaam), overdreven is, of niet ophoudt: dan
beschadigen haar \omterm{prop:g12:innate-immunity:inflammation}{mediatoren} en \omterm{def:g12:innate-immunity:phagocytosis}{fagocyten} het weefsel dat zij hadden moeten
beschermen. Wat de doorslag geeft is het doelwit en de duur, niet het
mechanisme, dat in beide gevallen hetzelfde is.
\end{solution}

\begin{solution}{pb:g12:innate-immunity:1}
\textbf{1.} Drie verdubbelingen in 2 uur: $\num{10000} \times 8 =
\num{80000}$; zes in 4 uur: $\num{640000}$.

\textbf{2.} $200 \times 5 = 1000$ per uur --- ver onder de
$\num{10000}$ nieuwe bacteriën die de verdubbelingen van het eerste uur
erbij doen: de groei gaat door.

\textbf{3.} $\num{50000} \times 20 = 10^6$ bacteriën.

\textbf{4.} Op uur 2 zijn er minder dan $\num{80000}$ bacteriën; de
neutrofielen van het eerste uur alleen al kunnen er een miljoen weghalen.
De \omterm{def:g12:innate-immunity:phagocytosis}{fagocyten} winnen vanaf uur 2.

\textbf{5.} Binnen enkele uren zijn de bacteriën weg; wat er overblijft
zijn de dode neutrofielen, het puin en het vocht die de pus zullen
vormen, en macrofagen die dat opruimen.

\textbf{6.} Roodheid en warmte: histamine en prostaglandinen verwijden de
vaten. Zwelling: histamine maakt de wanden lek en er sijpelt plasma uit.
Pijn: prostaglandinen maken de zenuwuiteinden gevoeliger.

\textbf{7.} De lekke vaten zijn hoe de neutrofielen en de verdedigende
\omterm{def:g10:chemistry-of-life:families}{eiwitten} van het plasma het weefsel bereiken; het vocht dat op de
gevoeliger gemaakte zenuwen drukt is de pijn.

\textbf{8.} Elke hartslag duwt meer bloed de verwijde vaten in; het
gezwollen weefsel kan niet uitzetten, dus verhoogt elke slag de druk op de
gevoeliger gemaakte zenuwen: een klop op de maat van het hart.

\textbf{9.} Dode neutrofielen, bacteriën dood en levend, verteerd puin,
plasma.

\textbf{10.} De patronen verdwenen met de bacteriën en het puin werd
opgeruimd: geen patroon, geen herkenning door de wachtcellen, geen
\omterm{prop:g12:innate-immunity:inflammation}{mediatoren}; de vaten worden weer normaal.

\textbf{11.} Het vermindert de vroege roodheid, warmte en lekkage die van
histamine komen; niet de latere zwelling en pijn die van prostaglandinen
komen, en ook niet het oproepen van de neutrofielen.

\textbf{12.} Het blokkeert het \omterm{def:g11:enzymes-and-phenotype:enzyme}{enzym} dat prostaglandinen maakt: de pijn,
een deel van de zwelling en de koorts nemen af.

\textbf{13.} De zalf zet de meeste \omterm{prop:g12:innate-immunity:inflammation}{mediatoren} stil, dus worden er minder
neutrofielen opgeroepen; met minder \omterm{def:g12:innate-immunity:phagocytosis}{fagocyten} worden de bacteriën
misschien niet opgeruimd en kan de infectie zich uitbreiden.

\textbf{14.} Alle drie werken op de tekenen en geen ervan doodt bacteriën;
de corticosteroïde kan, door het oproepen van \omterm{def:g12:innate-immunity:phagocytosis}{fagocyten} te knippen, de
infectie langer laten duren.

\textbf{15.} Cytokinen die de hersenen bereiken en de streefwaarde voor de
temperatuur verhogen; de extra warmte versnelt de \omterm{def:g12:innate-immunity:phagocytosis}{fagocyten} en vertraagt
de bacteriën.

\textbf{16.} De cytokinen die zij op de plek hebben opgenomen, die de
lymfocyten vertellen wat voor type indringer het was; zij brengen verslag
uit aan de lymfocyten van de lymfeknoop.

\textbf{17.} Ongeveer een week tot tien dagen. Voor een splinter die de
\omterm{def:g12:innate-immunity:immunity}{aangeboren afweer} in enkele uren opruimt zijn zij niet nodig --- maar zij
worden wel gemaakt, en onthouden.

\textbf{18.} De \omterm{def:g12:innate-immunity:immunity}{verworven afweer} is de tweede keer sneller en sterker,
dankzij het geheugen; de \omterm{def:g12:innate-immunity:immunity}{aangeboren afweer} verloopt precies hetzelfde.

\textbf{19.} \omterm{def:g11:antibiotic-resistance:antibiotic}{Antibiotica} doden bacteriën en nemen tegen bacteriële
infecties de plaats in van de ontbrekende \omterm{def:g12:innate-immunity:phagocytosis}{fagocyten}; zij kunnen niet het
opruimen van puin door de \omterm{def:g12:innate-immunity:phagocytosis}{fagocyten} vervangen, noch hun werking tegen
schimmels, noch het verslag aan de lymfeknoop.

\textbf{20.} Rond uur 2 streven de \omterm{def:g12:innate-immunity:phagocytosis}{fagocyten} de bacteriën voorbij;
histamine en prostaglandinen achter roodheid, warmte en zwelling,
prostaglandinen achter de pijn; de \omterm{prop:g12:innate-immunity:handover}{dendritische cel} draagt het verslag.
\end{solution}
